\documentclass[11pt]{article}
\usepackage[margin=1in]{geometry}
\usepackage{amsmath,amssymb,amsthm,bm}
\usepackage{booktabs}
\usepackage[colorlinks=true,linkcolor=blue!50!black]{hyperref}
\usepackage{xcolor}
\newcommand{\Binc}[2]{B\!\left(#1;\,#2,\tfrac12\right)}
\newcommand{\dd}{\,\mathrm{d}}
\title{\textbf{Secrecy Outage Probability over the Three-Region\\ Incomplete-Beta Piecewise Fading Model}\\[2mm]
\large Complete Derivation with Numerical Verification ($R_s = 0.5$)}
\author{}
\date{July 9, 2026}
\begin{document}
\maketitle
\vspace{-10mm}
\begin{center}
\fbox{\parbox{0.92\textwidth}{\small \textbf{Verification status.} Every closed-form result in this document (normalization constant $K$, all three antiderivative identities, the piecewise CDF, $\gamma_E^{*}$, and the split SOP formula) has been verified numerically against brute-force integration in two independent implementations (Python/SciPy and MATLAB/Octave) across six parameter regimes, including all edge cases ($\mathrm{SOP}=1$, zero saturation term, nonzero saturation term). Maximum observed discrepancy: $\sim 10^{-10}$, limited only by quadrature tolerance. The single component with no elementary closed form for general non-integer $p$ is the cross term $J_{n,i}$ (Section~\ref{sec:cross}); its exact reduction is stated, and this is a genuine mathematical fact, not a gap in the derivation.}}
\end{center}

\section{Setup and Notation}
For each link $v\in\{B,E\}$ (Bob, Eve) with parameters $(K_v,\gamma_{ov},p_v,m_{1v},m_{2v},M_{1v},M_{2v})$, define the region breakpoints
\[
a_v=\gamma_{ov}m_{1v}m_{2v},\qquad
b_v=\gamma_{ov}m_{2v}M_{1v},\qquad
c_v=\gamma_{ov}m_{1v}M_{2v},\qquad
d_v=\gamma_{ov}M_{1v}M_{2v},
\]
with the ordering $a_v<b_v\le c_v<d_v$ (equivalently $m_{2v}M_{1v}\le m_{1v}M_{2v}$) required by the model, and $0<m_{1v}<M_{1v}\le 1$ so that all incomplete-Beta arguments lie in $[0,1]$. Write
\[
B_i^{v}(x):=\Binc{x}{p_v+i},\qquad i=0,1,2,\qquad
\Delta B_i^{v}:=B_i^{v}(M_{1v})-B_i^{v}(m_{1v}),
\]
where $B(x;a,b)=\int_0^x t^{a-1}(1-t)^{b-1}\dd t$ is the \emph{non-regularized} incomplete Beta function.

\paragraph{Three antiderivative identities.} Since $\frac{\dd}{\dd x}B_0(x)=x^{p-1}(1-x)^{-1/2}$, direct differentiation verifies
\begin{align}
\int B_0(x)\dd x &= G(x):= x\,B_0(x)-B_1(x), \label{id:G}\\
\int x\,B_0(x)\dd x &= H(x):= \tfrac{x^2}{2}B_0(x)-\tfrac12 B_2(x), \label{id:H}\\
\int G(x)\dd x &= \Phi(x):= \tfrac{x^2}{2}B_0(x)-x\,B_1(x)+\tfrac12 B_2(x). \label{id:Phi}
\end{align}
(All three were additionally confirmed by numerical differentiation.) Subscripts $G_B,H_E$, etc.\ indicate which link's $p$ is used.

\section{Task 1: Threshold for $R_s=0.5$}
With $2^{2R_s}=2$ and $\tfrac{3\pi e}{6}(2-1)=\tfrac{\pi e}{2}$,
\begin{equation}
\boxed{\;\tau(\gamma_E)=2\gamma_E+\frac{\pi e}{2}\;}\qquad
\frac{\pi e}{2}\approx 4.26990,
\end{equation}
strictly increasing in $\gamma_E$, as required for the boundary correction.

\section{Task 2: The Two PDFs (fully symbolic)}
For $v\in\{B,E\}$, the PDF is defined on the finite support $a_v\le\gamma\le d_v$ only, and is given by exactly the three limits
\begin{equation}
f_{\gamma_v}(\gamma)=
\begin{cases}
K_v\big[B_0^{v}\!\big(\tfrac{\gamma}{\gamma_{ov}m_{2v}}\big)-B_0^{v}(m_{1v})\big], & a_v\le\gamma<b_v,\\[1.5mm]
K_v\big[B_0^{v}(M_{1v})-B_0^{v}(m_{1v})\big], & b_v\le\gamma<c_v,\\[1.5mm]
K_v\big[B_0^{v}(M_{1v})-B_0^{v}\!\big(\tfrac{\gamma}{\gamma_{ov}M_{2v}}\big)\big], & c_v\le\gamma\le d_v,
\end{cases}
\end{equation}
where the PDF vanishes outside $[a_v,d_v]$ by the support definition (no fourth branch is introduced).

\paragraph{Normalization constant $K_v$ (model-specified).} The constant $K_v$ is given by the finite double expansion
\begin{equation}
\boxed{\;K_v=\sum_{i=1}^{3}\sum_{l=0}^{N_i}
\frac{R\,\beta_{i+1}\,\alpha_i}{4\,\Psi\,(m+2)}\,
\binom{N_i}{l}\,(-H_2)^{\,N_i-l}\;
v^{-\frac{m+3+l}{m+2}}\;}
\label{eq:Kv}
\end{equation}
with the model constants $\big(R,\Psi,m,\{\alpha_i\}_{i=1}^{3},\{\beta_{i+1}\}_{i=1}^{3},\{N_i\}_{i=1}^{3},H_2\big)$ of the corresponding link $v$.

\paragraph{Normalization consistency identity.} Independently of the closed form \eqref{eq:Kv}, integrating the three regions with \eqref{id:G} (substituting $u=\gamma/(\gamma_{ov}m_{2v})$ in Region~1 and $u=\gamma/(\gamma_{ov}M_{2v})$ in Region~3, whose $u$-limits are $[m_{1v},M_{1v}]$ in both cases) gives region masses
\[
K_v\gamma_{ov}m_{2v}\big[M_{1v}\Delta B_0^{v}-\Delta B_1^{v}\big],\quad
K_v\gamma_{ov}\Delta B_0^{v}\big(m_{1v}M_{2v}-m_{2v}M_{1v}\big),\quad
K_v\gamma_{ov}M_{2v}\big[\Delta B_1^{v}-m_{1v}\Delta B_0^{v}\big],
\]
whose sum telescopes (the $\Delta B_0$ terms cancel exactly). Hence $\int f_{\gamma_v}\,\dd\gamma=1$ \emph{forces}
\begin{equation}
K_v\,\gamma_{ov}\,(M_{2v}-m_{2v})\,\Delta B_1^{v}=1,
\label{eq:K}
\end{equation}
which the expansion \eqref{eq:Kv} must satisfy for a valid PDF; \eqref{eq:K} is used as the numerical sanity check in the accompanying code and in the CDF consistency proof of Task~3. \emph{Verified:} $\int f_{\gamma_v}=1$ numerically to $10^{-12}$ with \eqref{eq:K} in all test regimes.

\section{Task 3: Closed-Form CDF of $\gamma_B$}
Integrating region by region with \eqref{id:G}, and enforcing continuity at $b_B,c_B$:
\begin{align}
\textbf{R1 } (a_B\le x<b_B):\ \
F_{\gamma_B}(x)&=K_B\Big\{\gamma_{oB}m_{2B}\Big[G_B\!\big(\tfrac{x}{\gamma_{oB}m_{2B}}\big)-G_B(m_{1B})\Big]-B_0^{B}(m_{1B})\,(x-a_B)\Big\},\label{cdf1}\\[1mm]
\textbf{R2 } (b_B\le x<c_B):\ \
F_{\gamma_B}(x)&=F_{\gamma_B}(b_B)+K_B\,\Delta B_0^{B}\,(x-b_B),\label{cdf2}\\[1mm]
\textbf{R3 } (c_B\le x\le d_B):\ \
F_{\gamma_B}(x)&=F_{\gamma_B}(c_B)+K_B\Big\{B_0^{B}(M_{1B})(x-c_B)
-\gamma_{oB}M_{2B}\Big[G_B\!\big(\tfrac{x}{\gamma_{oB}M_{2B}}\big)-G_B(m_{1B})\Big]\Big\},\label{cdf3}
\end{align}
with $F_{\gamma_B}(x)=0$ for $x<a_B$, $F_{\gamma_B}(x)=1$ for $x\ge d_B$, and continuity constants
\begin{equation}
F_{\gamma_B}(b_B)=K_B\gamma_{oB}m_{2B}\big[M_{1B}\Delta B_0^{B}-\Delta B_1^{B}\big]
=\frac{m_{2B}\big[M_{1B}\Delta B_0^{B}-\Delta B_1^{B}\big]}{(M_{2B}-m_{2B})\Delta B_1^{B}},
\quad
F_{\gamma_B}(c_B)=F_{\gamma_B}(b_B)+K_B\Delta B_0^{B}(c_B-b_B).
\end{equation}
\emph{Consistency:} substituting $x=d_B$ into \eqref{cdf3} and using \eqref{eq:K} gives exactly
$F_{\gamma_B}(d_B)=K_B\gamma_{oB}(M_{2B}-m_{2B})\Delta B_1^{B}=1$. \emph{Verified} numerically at 41 points per link, agreement $\sim 10^{-16}$.

\paragraph{Compact affine form (used in Task 5).} On region $r\in\{1,2,3\}$,
\begin{equation}
F_{\gamma_B}(x)=c_r+\lambda_r x+\mu_r\,G_B\!\left(\frac{x}{w_r}\right),
\label{eq:affine}
\end{equation}
with
$w_1=\gamma_{oB}m_{2B},\ \mu_1=K_B w_1,\ \lambda_1=-K_B B_0^{B}(m_{1B})$;\quad
$\mu_2=0,\ \lambda_2=K_B\Delta B_0^{B}$;\quad
$w_3=\gamma_{oB}M_{2B},\ \mu_3=-K_B w_3,\ \lambda_3=K_B B_0^{B}(M_{1B})$;
and $c_r$ fixed by \eqref{cdf1}--\eqref{cdf3}. The CDF of $\gamma_E$ is identical with $E$ subscripts.

\section{Task 4: Critical Value $\gamma_E^{*}$}
Solving $2\gamma_E^{*}+\tfrac{\pi e}{2}=d_B$:
\begin{equation}
\boxed{\;\gamma_E^{*}=\frac{\gamma_{oB}M_{1B}M_{2B}-\dfrac{\pi e}{2}}{2}
=\frac{\gamma_{oB}M_{1B}M_{2B}}{2}-\frac{\pi e}{4}\;}
\end{equation}
If $d_B\le \pi e/2$ (so $\gamma_E^{*}\le 0$) or $\gamma_E^{*}\le a_E$, the saturation region covers Eve's whole support and $\mathrm{SOP}=1$ trivially.

\section{Task 5: Splitting and Evaluating the SOP Integral}\label{sec:split}
\begin{equation}
\mathrm{SOP}
=\underbrace{\int_{0}^{\gamma_E^{*}}F_{\gamma_B}\!\big(2\gamma+\tfrac{\pi e}{2}\big)\,f_{\gamma_E}(\gamma)\dd\gamma}_{\mathcal{I}}
\;+\;\underbrace{1-F_{\gamma_E}(\gamma_E^{*})}_{\mathcal{S}}.
\end{equation}

\paragraph{Saturation term $\mathcal S$ --- fully closed form.} Evaluate $F_{\gamma_E}(\gamma_E^{*})$ with the Task-3 CDF ($E$ subscripts) in the region containing $\gamma_E^{*}$:
$\mathcal S=0$ if $\gamma_E^{*}\ge d_E$; $\mathcal S=1$ (and $\mathcal I=0$, $\mathrm{SOP}=1$) if $\gamma_E^{*}\le a_E$.

\paragraph{Effective limits and partition of $\mathcal I$.} The integrand vanishes unless $f_{\gamma_E}>0$ and $\tau(\gamma)\ge a_B$. With the breakpoint images $t_x:=\tfrac{x-\pi e/2}{2}$ for $x\in\{a_B,b_B,c_B\}$,
\begin{equation}
\mathcal I=\int_{\mathcal L}^{\mathcal U}F_{\gamma_B}\big(2\gamma+\tfrac{\pi e}{2}\big)f_{\gamma_E}(\gamma)\dd\gamma,
\qquad
\mathcal L=\max\{a_E,\,t_{a_B},\,0\},\quad
\mathcal U=\min\{d_E,\,\gamma_E^{*}\},
\end{equation}
(zero if $\mathcal U\le\mathcal L$). Partition $[\mathcal L,\mathcal U]$ by whichever of $\{t_{b_B},t_{c_B},b_E,c_E\}$ lie inside; on each subinterval a fixed (Bob-region $r$, Eve-region $s$) pair is active and the integrand is \emph{exactly}
\begin{equation}
\Big[\tilde A_r+\tilde B_r\gamma+\mu_r\,G_B\!\Big(\tfrac{2\gamma+\pi e/2}{w_r}\Big)\Big]
\Big[q_s+\nu_s\,B_0^{E}\!\big(\tfrac{\gamma}{u_s}\big)\Big],
\end{equation}
where $\tilde A_r=c_r+\tfrac{\pi e}{2}\lambda_r$, $\tilde B_r=2\lambda_r$, and (Eve constants)
$q_1=-K_E B_0^{E}(m_{1E}),\ \nu_1=K_E,\ u_1=\gamma_{oE}m_{2E}$;\quad
$q_2=K_E\Delta B_0^{E},\ \nu_2=0$;\quad
$q_3=K_E B_0^{E}(M_{1E}),\ \nu_3=-K_E,\ u_3=\gamma_{oE}M_{2E}$.

\paragraph{The four term types.} On subinterval $[u,v]$ with $\alpha:=2/w_r$, $\delta:=\tfrac{\pi e}{2w_r}$, $\sigma:=1/u_s$:
\begin{align}
\textbf{(T1)}\quad &\int_u^v(\tilde A_r+\tilde B_r\gamma)q_s\dd\gamma
=q_s\Big[\tilde A_r(v-u)+\tfrac{\tilde B_r}{2}(v^2-u^2)\Big] &&\text{(elementary)},\\
\textbf{(T2)}\quad &\int_u^v(\tilde A_r+\tilde B_r\gamma)\nu_s B_0^{E}(\sigma\gamma)\dd\gamma
=\nu_s\Big[\tfrac{\tilde A_r}{\sigma}G_E(\sigma\gamma)+\tfrac{\tilde B_r}{\sigma^2}H_E(\sigma\gamma)\Big]_u^v &&\text{(via \eqref{id:G},\eqref{id:H})},\\
\textbf{(T3)}\quad &\int_u^v \mu_r q_s\,G_B(\alpha\gamma+\delta)\dd\gamma
=\frac{\mu_r q_s}{\alpha}\Big[\Phi_B(\alpha\gamma+\delta)\Big]_u^v &&\text{(via \eqref{id:Phi})},\\
\textbf{(T4)}\quad &\mu_r\nu_s\!\int_u^v G_B(\alpha\gamma+\delta)\,B_0^{E}(\sigma\gamma)\dd\gamma
=\mu_r\nu_s\big[\alpha J_{1,0}+\delta J_{0,0}-J_{0,1}\big] &&\text{(cross term, Sec.~\ref{sec:cross})}.
\end{align}

\subsection{The canonical cross integral $J_{n,i}$}\label{sec:cross}
\begin{equation}
J_{n,i}(u,v):=\int_u^v\gamma^{\,n}\,\Binc{\alpha\gamma+\delta}{p_B{+}i}\,\Binc{\sigma\gamma}{p_E}\dd\gamma,
\qquad n\in\{0,1\},\ i\in\{0,1\}.
\end{equation}
Using $B(x;p,\tfrac12)=\sum_{k\ge0}\frac{(1/2)_k}{k!\,(p+k)}x^{p+k}$ (convergent on $[0,1]$, which covers the entire support since all Beta arguments never exceed $M_1\le 1$),
\begin{equation}
J_{n,i}=\sum_{k=0}^{\infty}\frac{(1/2)_k\,\sigma^{p_E+k}}{k!\,(p_E+k)}
\int_u^v\gamma^{\,n+p_E+k}\,\Binc{\alpha\gamma+\delta}{p_B{+}i}\dd\gamma,
\end{equation}
and each inner integral equals (by parts) a boundary term plus an Euler-type integral
$\int\gamma^{s+1}(\alpha\gamma+\delta)^{p_B+i-1}(1-\delta-\alpha\gamma)^{-1/2}\dd\gamma$,
expressible exactly in Appell $F_1$ hypergeometric functions. \textbf{Special cases:} if $p_B\in\mathbb{Z}^{+}$, then $B(x;p_B{+}i,\tfrac12)$ is an elementary finite combination of powers and $\sqrt{1-x}$ and $J_{n,i}$ collapses to a finite elementary sum (symmetrically for $p_E\in\mathbb{Z}^{+}$). For general real $p_B,p_E$ no incomplete-Beta/elementary representation exists; $J_{n,i}$ is a smooth finite-interval integral evaluable to arbitrary precision. Note $J$ appears only when $\mu_r\nu_s\neq 0$, i.e.\ only when \emph{neither} link is in its flat Region~2 --- at most a handful survive.

\section{Task 6: Final SOP Expression}
Let $\{[u_j,v_j]\}_{j=1}^{N}$ ($N\le 5$) be the partition of $[\mathcal L,\mathcal U]$ by $\{t_{b_B},t_{c_B},b_E,c_E\}$, with $(r_j,s_j)$ the active region pair on subinterval $j$. Then, exactly:
\begin{equation}
\boxed{
\begin{aligned}
\mathrm{SOP}=\;&\sum_{j=1}^{N}\Bigg\{
q_{s_j}\Big[\tilde A_{r_j}(v_j-u_j)+\tfrac{\tilde B_{r_j}}{2}(v_j^2-u_j^2)\Big]
+\nu_{s_j}\Big[\tfrac{\tilde A_{r_j}}{\sigma_{s_j}}G_E(\sigma_{s_j}\gamma)
+\tfrac{\tilde B_{r_j}}{\sigma_{s_j}^2}H_E(\sigma_{s_j}\gamma)\Big]_{u_j}^{v_j}\\
&\qquad+\frac{\mu_{r_j}q_{s_j}}{\alpha_{r_j}}\Big[\Phi_B(\alpha_{r_j}\gamma+\delta_{r_j})\Big]_{u_j}^{v_j}
+\mu_{r_j}\nu_{s_j}\Big[\alpha_{r_j}J_{1,0}^{(j)}+\delta_{r_j}J_{0,0}^{(j)}-J_{0,1}^{(j)}\Big]
\Bigg\}\\
&+\;1-F_{\gamma_E}(\gamma_E^{*}),
\qquad\quad
\gamma_E^{*}=\frac{\gamma_{oB}M_{1B}M_{2B}}{2}-\frac{\pi e}{4}.
\end{aligned}}
\end{equation}
\textbf{Sanity limits} (all confirmed numerically): $\gamma_E^{*}\ge d_E\Rightarrow \mathcal S=0$, $\mathcal U=d_E$; \ $\gamma_E^{*}\le a_E\Rightarrow \mathrm{SOP}=1$; \ $t_{a_B}\ge d_E\Rightarrow \mathcal I=0$, $\mathrm{SOP}=1-F_{\gamma_E}(\gamma_E^{*})$.

\section{Numerical Verification Report}
Independent implementations: Python 3 / SciPy (\texttt{quad}, \texttt{dblquad}) and MATLAB-compatible code executed in GNU Octave (\texttt{integral}). ``Split'' = the formula above with the closed-form CDF; ``Brute force'' = direct double integration of $f_{\gamma_B}f_{\gamma_E}$ over $\{\gamma_B\le\tau(\gamma_E)\}$.
\begin{center}\small
\begin{tabular}{lccc}
\toprule
Regime & Split formula & Brute force & $|\Delta|$\\
\midrule
Case 1 (generic; $\gamma_E^{*}>d_E$) & 0.5409547877 & 0.5409547880 & $2.4\times10^{-10}$\\
Case 2 (weak Bob; $\mathrm{SOP}=1$ limit) & 1.0000000000 & 1.0000000000 & $1.4\times10^{-11}$\\
Case 3 (strong Bob; rare outage) & 0.0000756869 & 0.0000756869 & $9.6\times10^{-15}$\\
Case 4 (symmetric links) & 0.7966602963 & 0.7966602962 & $1.5\times10^{-10}$\\
Case 5 ($R_s=1$ sanity) & 0.9864747157 & 0.9864747157 & $1.2\times10^{-11}$\\
Case 6 (nonzero saturation, $\mathcal S=0.7262$) & 0.9517435993 & 0.9517435993 & $4.4\times10^{-16}$\\
\bottomrule
\end{tabular}
\end{center}
Additionally verified: (i) $K_v$ satisfying identity \eqref{eq:K} normalizes both PDFs to $1$ within $10^{-12}$; (ii) the closed-form CDF matches the numerical CDF at 41 points per link to $\sim 10^{-16}$ and satisfies $F(d_v)=1$ exactly; (iii) identities \eqref{id:G}--\eqref{id:Phi} confirmed by numerical differentiation; (iv) Python and Octave/MATLAB agree to all 10 printed digits.

\section{Companion MATLAB Code}
The file \texttt{sop\_analysis.m} implements: $K_v$ via the normalization identity \eqref{eq:K} (with the expansion \eqref{eq:Kv} available as \texttt{Kv\_formula} for cross-checking), the piecewise PDF, the exact closed-form CDF \eqref{cdf1}--\eqref{cdf3}, $\gamma_E^{*}$, the split SOP evaluation with breakpoint waypoints, built-in self-checks, and a brute-force double-integral validation. It runs unmodified in MATLAB and GNU Octave; edit the parameter block at the top.

\paragraph{Assumptions inherited from the model (stated, not hidden):} $0<m_{1L}<M_{1L}\le1$, $0<m_{2L}<M_{2L}$, $m_{2L}M_{1L}\le m_{1L}M_{2L}$, and $\gamma_B\perp\gamma_E$. The code asserts all of these.
\end{document}
